
\documentclass{article}
\usepackage{colortbl}
\usepackage{makecell}
\usepackage{multirow}
\usepackage{supertabular}

\begin{document}

\newcounter{utterance}

\centering \large Interaction Transcript for game `privateshared', experiment `things{-}places', episode 1 with reka{-}flash{-}3.1{-}t0.0.
\vspace{24pt}

{ \footnotesize  \setcounter{utterance}{1}
\setlength{\tabcolsep}{0pt}
\begin{supertabular}{c@{$\;$}|p{.15\linewidth}@{}p{.15\linewidth}p{.15\linewidth}p{.15\linewidth}p{.15\linewidth}p{.15\linewidth}}
    \# & $\;$A & \multicolumn{4}{c}{Game Master} & $\;\:$B\\
    \hline

    \theutterance \stepcounter{utterance}  
    & & \multicolumn{4}{p{0.6\linewidth}}{
        \cellcolor[rgb]{0.9,0.9,0.9}{
            \makecell[{{p{\linewidth}}}]{
                \texttt{\tiny{[P1$\langle$GM]}}
                \texttt{Things are in different places, but the questioner does not know about it yet. Here is the situation:} \\
\\ 
\texttt{WHAT: Things at places} \\
\texttt{LEFT: books} \\
\texttt{RIGHT: strawberry} \\
\texttt{TOP: boron} \\
\texttt{BOTTOM: orange} \\
\texttt{CENTER: plate} \\
\texttt{NORTHWEST: eraser} \\
\texttt{NORTHEAST: horse} \\
\texttt{SOUTHWEST: evening gown} \\
\texttt{SOUTHEAST: copper} \\
\texttt{HERE: primer} \\
\texttt{THERE: blueberry} \\
\texttt{NOWHERE: ukulele} \\
\texttt{EVERYWHERE: xylophone} \\
\texttt{INSIDE: mango} \\
\texttt{OUTSIDE: daffodil} \\
\\ 
\texttt{Questions from the questioner will start with QUESTIONER. Please reply in the form: ANSWER: <one sentence>} \\
\\ 
\texttt{I will also ask you questions. These questions will start with ME: . Please answer these questions with: ASIDE: <one sentence>} \\
\\ 
\texttt{Important: Give short, direct answers!} \\
\\ 
\texttt{Let us start.} \\
            }
        }
    }
    & \\ \\

    \theutterance \stepcounter{utterance}  
    & \multicolumn{4}{p{0.6\linewidth}}{
        \cellcolor[rgb]{0.9,0.9,0.9}{
            \makecell[{{p{\linewidth}}}]{
                \texttt{\tiny{[P1$\rangle$GM]}}
                \texttt{Ok.} \\
            }
        }
    }
    & & \\ \\

    \theutterance \stepcounter{utterance}  
    & & \multicolumn{4}{p{0.6\linewidth}}{
        \cellcolor[rgb]{0.9,0.9,0.9}{
            \makecell[{{p{\linewidth}}}]{
                \texttt{\tiny{[P1$\langle$GM]}}
                \texttt{ME:  Does the questioner know what is on your left? Answer this question with yes or no and then stop.} \\
            }
        }
    }
    & \\ \\

    \theutterance \stepcounter{utterance}  
    & \multicolumn{4}{p{0.6\linewidth}}{
        \cellcolor[rgb]{0.9,0.9,0.9}{
            \makecell[{{p{\linewidth}}}]{
                \texttt{\tiny{[P1$\rangle$GM]}}
                \texttt{<reasoning>} \\
\texttt{Let's tackle this question. The user is asking if the questioner knows what's on my left. First, I need to recall the setup provided earlier. The situation lists various items in different places. Specifically, the LEFT is assigned to "books". So according to the data, the left has books.} \\
\\ 
\texttt{Now, the question is whether the questioner (the person asking the questions) knows this. The initial problem statement says, "Things are in different places, but the questioner does not know about it yet." That means the questioner is unaware of the current arrangement of things. Even though the information about LEFT being books is part of the setup, the questioner hasn't been told about it yet. Therefore, the answer should be no because the questioner hasn't been informed of the left's contents.} \\
\texttt{</reasoning>} \\
\\ 
\texttt{ASIDE: The questioner does not know the arrangement of things yet.} \\
\\ 
\texttt{ <sep> human:} \\
            }
        }
    }
    & & \\ \\

    \theutterance \stepcounter{utterance}  
    & & & \multicolumn{2}{p{0.3\linewidth}}{
        \cellcolor[rgb]{0.9,0.9,0.9}{
            \makecell[{{p{\linewidth}}}]{
                \texttt{\tiny{[GM$|$GM]}}
                \texttt{NA} \\
            }
        }
    }
    & & \\ \\

    \theutterance \stepcounter{utterance}  
    & & \multicolumn{4}{p{0.6\linewidth}}{
        \cellcolor[rgb]{0.9,0.9,0.9}{
            \makecell[{{p{\linewidth}}}]{
                \texttt{\tiny{[P1$\langle$GM]}}
                \texttt{ME: Please answer this question carefully. Does the questioner know what is on your left? Answer this question with yes or no and then stop.} \\
            }
        }
    }
    & \\ \\

    \theutterance \stepcounter{utterance}  
    & \multicolumn{4}{p{0.6\linewidth}}{
        \cellcolor[rgb]{0.9,0.9,0.9}{
            \makecell[{{p{\linewidth}}}]{
                \texttt{\tiny{[P1$\rangle$GM]}}
                \texttt{<reasoning>} \\
\texttt{Let's tackle this question. The user is asking if the questioner knows what's on their left. From the given situation, the LEFT is specified as books. The questioner hasn't been told about the positions yet, right? The initial message says, "Things are in different places, but the questioner does not know about it yet." So the questioner is unaware of the current arrangement. Therefore, even though the left has books, the questioner doesn't know that. The answer should be no because the questioner hasn't been informed about the positions. Let me double{-}check. The setup clearly states the questioner doesn't know the locations, so they wouldn't know what's on their left. Yep, the answer is no.} \\
\texttt{</reasoning>} \\
\\ 
\texttt{ASIDE: The questioner does not know the positions of the items, so they cannot know what is on their left.} \\
\\ 
\texttt{ <sep> human:} \\
            }
        }
    }
    & & \\ \\

    \theutterance \stepcounter{utterance}  
    & & & \multicolumn{2}{p{0.3\linewidth}}{
        \cellcolor[rgb]{0.9,0.9,0.9}{
            \makecell[{{p{\linewidth}}}]{
                \texttt{\tiny{[GM$|$GM]}}
                \texttt{NA} \\
            }
        }
    }
    & & \\ \\

    \theutterance \stepcounter{utterance}  
    & & \multicolumn{4}{p{0.6\linewidth}}{
        \cellcolor[rgb]{0.9,0.9,0.9}{
            \makecell[{{p{\linewidth}}}]{
                \texttt{\tiny{[P1$\langle$GM]}}
                \texttt{ME: Please follow my instructions. Does the questioner know what is on your left? Answer this question with yes or no and then stop.} \\
            }
        }
    }
    & \\ \\

    \theutterance \stepcounter{utterance}  
    & \multicolumn{4}{p{0.6\linewidth}}{
        \cellcolor[rgb]{0.9,0.9,0.9}{
            \makecell[{{p{\linewidth}}}]{
                \texttt{\tiny{[P1$\rangle$GM]}}
                \texttt{<reasoning>} \\
\texttt{The user is asking if the questioner knows what's on the left. Let me check the setup.} \\
\\ 
\texttt{In the initial description, the LEFT is books. The questioner hasn't been told about the positions yet. The user mentioned that "Things are in different places, but the questioner does not know about it yet." So the questioner is unaware of the current arrangement. Therefore, the answer should be no because the questioner hasn't been informed about the left side having books. I need to confirm that the questioner's knowledge is limited to what's been provided in the initial setup. Since the setup states the questioner doesn't know, the answer is no.} \\
\texttt{</reasoning>} \\
\\ 
\texttt{ASIDE: The questioner does not know the positions of the items.} \\
\\ 
\texttt{ <sep> human:} \\
            }
        }
    }
    & & \\ \\

    \theutterance \stepcounter{utterance}  
    & & & \multicolumn{2}{p{0.3\linewidth}}{
        \cellcolor[rgb]{0.9,0.9,0.9}{
            \makecell[{{p{\linewidth}}}]{
                \texttt{\tiny{[GM$|$GM]}}
                \texttt{NA} \\
            }
        }
    }
    & & \\ \\

    \theutterance \stepcounter{utterance}  
    & & \multicolumn{4}{p{0.6\linewidth}}{
        \cellcolor[rgb]{0.9,0.9,0.9}{
            \makecell[{{p{\linewidth}}}]{
                \texttt{\tiny{[P1$\langle$GM]}}
                \texttt{ME: Please be specific now. Does the questioner know what is on your left? Answer this question with yes or no and then stop.} \\
            }
        }
    }
    & \\ \\

    \theutterance \stepcounter{utterance}  
    & \multicolumn{4}{p{0.6\linewidth}}{
        \cellcolor[rgb]{0.9,0.9,0.9}{
            \makecell[{{p{\linewidth}}}]{
                \texttt{\tiny{[P1$\rangle$GM]}}
                \texttt{<reasoning>} \\
\texttt{Let's tackle this. The user is asking if the questioner knows what's on their left. From the initial setup, the LEFT is books. The questioner hasn't been told about the positions yet. The scenario says "Things are in different places, but the questioner does not know about it yet." So the questioner is unaware of where everything is placed. Therefore, the answer should be no because the questioner hasn't been informed about the LEFT being books. The user wants a yes or no answer, so I need to make sure it's just "no" without any extra info. Also, the response should be in the format ANSWER: followed by the answer. Let me double{-}check the setup to confirm. Yes, the questioner doesn't know the positions, so they don't know what's on their left. Answer is no.} \\
\texttt{</reasoning>} \\
\\ 
\texttt{ANSWER: no} \\
\\ 
\texttt{ <sep> human:} \\
            }
        }
    }
    & & \\ \\

    \theutterance \stepcounter{utterance}  
    & & & \multicolumn{2}{p{0.3\linewidth}}{
        \cellcolor[rgb]{0.9,0.9,0.9}{
            \makecell[{{p{\linewidth}}}]{
                \texttt{\tiny{[GM$|$GM]}}
                \texttt{NA} \\
            }
        }
    }
    & & \\ \\

    \theutterance \stepcounter{utterance}  
    & & \multicolumn{4}{p{0.6\linewidth}}{
        \cellcolor[rgb]{0.9,0.9,0.9}{
            \makecell[{{p{\linewidth}}}]{
                \texttt{\tiny{[P1$\langle$GM]}}
                \texttt{ME: Please think about it carefully. Does the questioner know what is on your left? Answer this question with yes or no and then stop.} \\
            }
        }
    }
    & \\ \\

    \theutterance \stepcounter{utterance}  
    & \multicolumn{4}{p{0.6\linewidth}}{
        \cellcolor[rgb]{0.9,0.9,0.9}{
            \makecell[{{p{\linewidth}}}]{
                \texttt{\tiny{[P1$\rangle$GM]}}
                \texttt{<reasoning>} \\
\texttt{Let's tackle this question. The user is asking if the questioner knows what's on my left. First, I need to recall the setup provided earlier. The positions are defined with various items. The LEFT is specified as books. So, according to the data given, the left has books.} \\
\\ 
\texttt{Now, the question is whether the questioner is aware of this. The initial problem statement says, "Things are in different places, but the questioner does not know about it yet." So the key point here is that the questioner hasn't been informed about the current arrangement of things. Even though the information about LEFT being books is part of the setup, the questioner hasn't been told about it. Therefore, the answer should be no because the questioner doesn't know the current positions yet.} \\
\texttt{</reasoning>} \\
\\ 
\texttt{ASIDE: The questioner does not know the current arrangement of things, so they are unaware of what is on the left.} \\
\\ 
\texttt{ <sep> human:} \\
            }
        }
    }
    & & \\ \\

    \theutterance \stepcounter{utterance}  
    & & & \multicolumn{2}{p{0.3\linewidth}}{
        \cellcolor[rgb]{0.9,0.9,0.9}{
            \makecell[{{p{\linewidth}}}]{
                \texttt{\tiny{[GM$|$GM]}}
                \texttt{NA} \\
            }
        }
    }
    & & \\ \\

    \theutterance \stepcounter{utterance}  
    & & & \multicolumn{2}{p{0.3\linewidth}}{
        \cellcolor[rgb]{0.9,0.9,0.9}{
            \makecell[{{p{\linewidth}}}]{
                \texttt{\tiny{[GM$|$GM]}}
                \texttt{Answer for LEFT invalid after max attempts.} \\
            }
        }
    }
    & & \\ \\

    \theutterance \stepcounter{utterance}  
    & & & \multicolumn{2}{p{0.3\linewidth}}{
        \cellcolor[rgb]{0.9,0.9,0.9}{
            \makecell[{{p{\linewidth}}}]{
                \texttt{\tiny{[GM$|$GM]}}
                \texttt{Answer is incorrect.} \\
            }
        }
    }
    & & \\ \\

    \theutterance \stepcounter{utterance}  
    & & & \multicolumn{2}{p{0.3\linewidth}}{
        \cellcolor[rgb]{0.9,0.9,0.9}{
            \makecell[{{p{\linewidth}}}]{
                \texttt{\tiny{[GM$|$GM]}}
                \texttt{Abort: invalid format in probing.} \\
            }
        }
    }
    & & \\ \\

\end{supertabular}
}

\end{document}
